\chapter{Le sondage : résultats}\label{chap:04-sondage-resultats}

Ce chapitre donne tous les résultats du sondage : les échantillons obtenus
(section~\ref{sec:sr-echantillons}), la comparaison principale à grille égale, par critère
(section~\ref{sec:sr-principale}) puis par faille (section~\ref{sec:sr-failles}), ce qui résiste aux
analyses de robustesse (section~\ref{sec:sr-robustesse}), notre lecture de 2018 contre les chiffres de
Pineau (section~\ref{sec:sr-pineau}), le premier tirage de 2026 et les articles sur les modèles de
langue (section~\ref{sec:sr-premier}), la trajectoire de 2019 à 2024 (section~\ref{sec:sr-trajectoire}),
l'accord entre lecteurs (section~\ref{sec:sr-accord}), enfin les analyses A5 et A6
(sections~\ref{sec:sr-declarations} et~\ref{sec:sr-filtre}). Tous les nombres viennent des fichiers
JSON de \file{study/} nommés dans chaque légende.

\section{Les échantillons obtenus}\label{sec:sr-echantillons}

\paragraph{Le premier tirage de 2026}
Les 50 articles tirés parmi les \num{1252} articles de RL de 2025 et 2026 se répartissent ainsi : 27
appliquent ce qu'ils appellent RL à de grands modèles de langue (famille \file{llm}), 14 relèvent du RL
classique (\file{classic}) et 9 sont théoriques (\file{theory}), dont 5 sans aucune expérience. Les
pourcentages se calculent sur les 45 articles qui ont des expériences.

\paragraph{Notre lecture de 2018}
Les 50 articles de 2018 (21 d'ICML, 20 de NeurIPS, 9 d'ICLR) comptent 48 articles expérimentaux. Sur
ces 48 articles, un article (C14) compare une méthode exacte sans hyperparamètre : Q3 et Q6 y valent
\file{na}, d'où un dénominateur de 47 pour ces deux critères. Q4 et Y5 valent \file{na} quand un
article ne compare aucune référence.

\paragraph{Le RL classique de 2026}
Les 50 articles réunissent les 14 articles \file{classic} du premier tirage et les 36 premiers articles
\file{classic} de la file A3 (\file{extension\_stats.json}, champ \file{samples}). Les 50 ont des
expériences. Atteindre le trente-sixième article \file{classic} de la file a demandé 57 lectures, dont 18 articles
\file{llm} et 3 articles \file{theory} lus en route malgré le tri sur le titre. Les 13 articles lus
après lui portent le total à 70 lectures.

\paragraph{\enquote{50 contre 48}}
L'article de blog et nos figures parlent de \enquote{50 articles de RL classique de 2026} contre
\enquote{50 articles de 2018}. Les deux échantillons comptent bien 50 articles lus. Les pourcentages de
2018 se calculent pourtant sur 48 articles, parfois 47, parce que deux articles de 2018 n'ont aucune
expérience. Ceux de 2026 se calculent sur 50 articles, parfois 49 quand un critère ne s'applique pas.
Le dénominateur exact de chaque part figure dans les tableaux de ce chapitre. Le libellé \enquote{50
contre 50} décrit donc les échantillons tirés, le libellé \enquote{48 contre 50} les échantillons
analysés pour la plupart des critères.

\section{La comparaison principale, à grille égale}\label{sec:sr-principale}

\paragraph{Les deux échantillons et leurs lecteurs}
La comparaison principale oppose nos 50 articles de RL classique de 2026 à notre lecture de 50 articles
de 2018, avec les mêmes grilles et les mêmes règles harmonisées (section~\ref{sec:sm-verdict}). Les
articles de 2018 et ceux de 2026 ont été lus par des lecteurs distincts (\file{study/survey\_extra.json},
champ \file{readers\_year}). Un effet propre aux lecteurs d'une année ne se sépare donc pas de l'année
dans cette comparaison (section~\ref{sec:sm-annees}). L'écart 36 les sépare en faisant relire les 98
articles par huit lecteurs qui lisent chacun les deux années (section~\ref{sec:sr-annees}). Le critère
X8, plus bas, échappait déjà à cette confusion.

Le tableau~\ref{tab:sondage-pineau} donne les critères de la grille v1, le
tableau~\ref{tab:sondage-extension} ceux de la grille v2. Pour chaque critère, ils donnent la part en
2018 et en 2026 avec leurs intervalles de Wilson, la différence avec son intervalle de Newcombe, la
valeur $p$ du test exact de Fisher et la valeur ajustée par la correction de Holm, appliquée séparément
aux 11 critères v1 et aux 12 critères v2. La figure~\ref{fig:forest-criteres} dessine ces différences
critère par critère. La figure~\ref{fig:grille-egale} reprend les huit pratiques de l'article de blog.

\gentable{sondage-pineau}

\gentable{sondage-extension}

\fig{forest-criteres}{Tous les critères de lecture, grille v1 (Q, X) et grille v2 (Y), dans notre
lecture de 2018 (48 articles expérimentaux) et dans nos 50 articles de RL classique de 2026, lus avec
les mêmes grilles et les mêmes règles. Pour chaque critère, les deux points donnent les deux parts,
la barre l'intervalle de Newcombe à 95~\% de leur différence. Les critères dont la différence passe la correction
de Holm sont mis en évidence. Pour X3, une part plus basse est un progrès. Données :
\file{study/extension\_stats.json}.}

\fig{grille-egale}{Différence, en points de pourcentage, entre 50 articles de RL classique de 2026 et 50
articles de 2018 (48 expérimentaux), lus avec les mêmes grilles et les mêmes règles, pour les huit
pratiques de l'article de blog. Le point donne la différence observée, la barre son intervalle de Newcombe à
95~\%. Les quatre pratiques en vert passent la correction de Holm : le lien vers le code (Q5),
l'effet d'un réglage montré (Y9), le calcul chiffré (X7), une incertitude rapportée (X2). Les quatre
autres, dont le test statistique (Q8), ont des intervalles qui contiennent zéro. Données :
\file{study/extension\_stats.json}.}

\paragraph{Quatre pratiques ont changé}
Quatre différences passent la correction de Holm.
\begin{itemize}
  \item \textbf{Le code en lien (Q5)} passe de 31~\% (15 sur 48) à 74~\% (37 sur 50) : $+43$ points,
  $[+23\,;\,+58]$, $p_{\text{Holm}} = 0{,}0004$.
  \item \textbf{La sensibilité à un hyperparamètre montrée (Y9)} passe de 44~\% (21 sur 48) à 80~\% (40
  sur 50) : $+36$ points, $[+17\,;\,+52]$, $p_{\text{Holm}} = 0{,}004$.
  \item \textbf{Le calcul chiffré (X7)} passe de 27~\% (13 sur 48) à 58~\% (29 sur 50) : $+31$ points,
  $[+11\,;\,+47]$, $p_{\text{Holm}} = 0{,}02$.
  \item \textbf{Une incertitude rapportée (X2)} passe de 63~\% (30 sur 48) à 88~\% (44 sur 50) : $+26$
  points, $[+8\,;\,+41]$, $p_{\text{Holm}} = 0{,}04$. C'est le verdict le plus fragile
  (section~\ref{sec:sr-robustesse}). Il compte toute barre d'erreur, quelle que soit la variabilité qu'elle
  mesure : le critère X8, plus bas, demande laquelle.
\end{itemize}
Un cinquième critère passe aussi la correction : l'emploi de réseaux de neurones (Q2), de 79~\% (38
sur 48) à 98~\% (49 sur 50), $+19$ points, $[+6\,;\,+32]$, $p_{\text{Holm}} = 0{,}03$. Q2 décrit la
population : le RL classique de 2026 est presque entièrement du RL profond, ce qu'il n'était pas tout à
fait en 2018. Le critère ne mesure aucune faille de reproductibilité.

\paragraph{Une barre d'erreur sur quoi ? Le critère X8}
X2 vaut \file{yes} dès qu'un article rapporte une incertitude. Or l'expérience B1 (\cref{sec:b1}) montre qu'une
barre calculée sur les épisodes d'un même entraînement est le piège même qui produit le $t$ de 9 de 2018 : elle
ignore la variabilité entre graines aléatoires, celle qui décide si deux algorithmes diffèrent (faille F11). Un
article qui trace ce genre de barre compte pourtant \file{yes} pour X2. L'écart 32, écrit avant toute lecture,
ajoute donc un critère exploratoire, X8 : l'incertitude rapportée est-elle calculée sur des graines aléatoires,
c'est-à-dire sur des entraînements indépendants ? Six agents IA neufs, qui n'avaient lu aucun article du
sondage, se sont partagé les 74 articles qui valent \file{yes} pour X2. Deux autres en ont relu 20 à
l'aveugle. Chacun des six lisait des articles des deux années, ce qui défait, pour ce seul critère, la confusion entre lecteurs et année. Parmi les articles expérimentaux, X8
passe de 35~\% (17 sur 48) à 42~\% (21 sur 50) : $+7$ points, $[-12\,;\,+25]$, $p = 0{,}54$
(\cref{tab:x8,fig:x2-x8}). La hausse de X2 ne s'accompagne d'aucune hausse détectable de l'incertitude
calculée sur des graines aléatoires. Les barres d'erreur se sont répandues sans que leur source change de façon
visible. En 2026, 14 des 44 articles qui tracent une barre ne disent pas sur quoi elle est calculée. X8 a été
décidé après les résultats principaux et lu par d'autres agents que les lecteurs principaux : il nuance le
verdict de X2 sans le remplacer.

\fig{x2-x8}{Un carré par article expérimental, 48 en 2018 et 50 en 2026, coloré selon ce sur quoi son
incertitude est calculée : graines aléatoires, graines et épisodes mêlés, épisodes, éléments de test, autre
source, source non dite, aucune incertitude. Les articles qui rapportent une incertitude passent de 63~\% à
88~\% (X2, différence qui passe la correction de Holm), ceux qui la calculent sur des graines aléatoires de 35~\% à
42~\% seulement (X8, $+7$ points, $[-12\,;\,+25]$, aucun changement détectable). En 2026, 14 des 44 articles
qui tracent une barre d'erreur ne disent pas sur quoi elle est calculée. Données : \file{study/x8\_stats.json},
\file{study/extension\_stats.json} (\cref{tab:x8}).}

\gentable{x8}

\paragraph{Deux intervalles qui évitent zéro sans passer la correction}
La façon de choisir les hyperparamètres (Q6) passe de 26~\% (12 sur 47) à 48~\% (24 sur 50) : $+22$
points, $[+3\,;\,+39]$, $p = 0{,}03$ avant correction, $p_{\text{Holm}} = 0{,}24$. Les outils
statistiques postérieurs à 2018 (Y12) passent de 2~\% (1 sur 48) à 14~\% (7 sur 50) : $+12$ points,
$[+1\,;\,+24]$, $p_{\text{Holm}} = 0{,}66$. Pris seuls, ces deux intervalles évitent zéro. Avec 11 et
12 critères testés à la fois, ils ne passent pas le seuil ajusté. Nous ne détectons donc aucun
changement après correction, sans pouvoir en exclure un.

\paragraph{Aucun changement détectable pour les autres critères}
Les 16 autres critères ont des intervalles qui contiennent zéro. Ceux qui mesurent une faille se rangent
en trois groupes.
\begin{itemize}
  \item \textbf{Des hausses que le sondage ne peut ni établir ni exclure.} Les hyperparamètres de la méthode
  (Q3, 55~\% puis 72~\%), l'évaluation tenue à part (Q7, 33~\% puis 45~\%) et le test de généralisation (Y8,
  23~\% puis 34~\%) montent de plus de dix points, comme le choix des hyperparamètres (Q6).
  \item \textbf{Des critères à plat.} Les hyperparamètres des références (Q4, 35~\% puis 29~\%), les
  références réglées avec le même soin (X4, 13~\% puis 16~\%), la mesure définie (Y1, 79~\% puis 78~\%), les
  meilleurs essais gardés (X3, 19~\% puis 14~\%), l'origine du code des références (Y5, 46~\% puis 41~\%) et
  la version des bancs d'essai (Y7, 23~\% puis 26~\%) bougent de sept points au plus. Le réglage des
  références, en particulier, ne montre aucun signe de progrès.
  \item \textbf{Des pratiques rares.} Le test statistique (Q8, 6~\% puis 12~\%) et les versions des
  bibliothèques (Y10, 2~\% puis 8~\%) partent de si bas qu'une hausse de six points reste dans le bruit du
  tirage.
\end{itemize}

\paragraph{La puissance du sondage}
Chaque intervalle de Newcombe s'étend sur vingt à quarante points. L'écart 30 (b), écrit avant le calcul
du dépôt, chiffre ce que le sondage pouvait voir (\cref{tab:difference-detectable}). Partant de la part de 2018 et des
effectifs réels, le test exact de Fisher bilatéral ne détecte avec 80~\% de chances qu'une hausse de 18 à 30
points selon le critère au seuil de 5~\%, de 26 à 39 points au seuil de Bonferroni sur onze critères. Pour Q7,
il faut 29 points, 37 après correction. La hausse effectivement vue, 12 points, n'avait que 18~\% de chances
d'être détectée au seuil de 5~\%, 4~\% après correction. Les quatre changements établis dépassent tous leur
différence détectable. X2 la dépasse de peu : la puissance du test pour la différence vue tombe à 52~\% au seuil de
Bonferroni. \enquote{Aucun changement détectable} ne veut donc pas dire \enquote{aucun changement}. Pour Q3 par
exemple, la différence observée vaut $+17$ points, avec un intervalle de $-2$ à $+34$ : la donnée reste compatible
avec une hausse de trente points comme avec une stagnation.

\gentable{difference-detectable}

\section{Relecture faille par faille}\label{sec:sr-failles}

Cette section reprend les résultats dans l'ordre des failles du tableau~\ref{tab:failles}. Les chiffres de 2018
et de 2026 sont ceux de la comparaison principale, sauf mention du premier tirage.

\paragraph{F1, le code}
Le lien vers le code passe de 31~\% à 74~\% des articles. La faille est écoutée. Le bilan officiel du
programme de reproductibilité de NeurIPS 2019 va dans le même sens : la part d'articles acceptés avec
du code passe de moins de 50~\% à 74,4~\% \parencite{pineau2021improving}. Dans tout le premier tirage,
grands modèles de langue compris, la part vaut 67~\% (30 sur 45). Un article de RL classique sur quatre
publie encore son résultat sans son code. La hausse se retrouve quand les deux années sont lues sur la même version
de l'article (section~\ref{sec:sr-robustesse}).

\paragraph{F2 et F3, le plan pour refaire et le matériel}
Le calcul est chiffré dans 58~\% des articles de 2026 contre 27~\% en 2018. Les hyperparamètres de la
méthode sont tous donnés dans 72~\% des articles contre 55~\%, sans changement détectable
($p_{\text{Holm}} = 0{,}58$). Les versions des bibliothèques restent rares : 4 articles de RL classique
sur 50 les donnent en 2026, contre 1 sur 48 en 2018 (Y10). Dans tout le premier tirage, un seul article
sur 45 les donne. Les deux failles sont écoutées en partie : le calcul se chiffre, les versions des
logiciels se taisent.

\paragraph{F4 et F5, les hyperparamètres de la méthode et leur choix}
Le niveau de Q3 ne bouge pas de façon détectable. La façon de choisir les hyperparamètres est décrite
dans 48~\% des articles de 2026 contre 26~\% en 2018, une différence qui ne passe pas la correction. Moins
d'un article sur deux dit encore comment il a trouvé ses réglages.

\paragraph{F6, la sensibilité aux réglages}
Au moins un hyperparamètre de la méthode est balayé, résultat à l'appui, dans 80~\% des articles de
2026 contre 44~\% en 2018. Cette faille aussi a été écoutée. Le premier tirage, grands modèles compris, donne un
niveau voisin : 32 articles sur 45.

\paragraph{F7, les références}
Les hyperparamètres des références sont tous donnés dans 29~\% des articles de 2026 (14 sur 49), contre
35~\% en 2018 (17 sur 48). Les articles qui disent avoir réglé leurs références avec le même soin que
leur méthode passent de 13~\% à 16~\%. La pratique reste rare les deux années, sans changement
détectable. Contre les 60~\% de Pineau, la baisse paraissait nette (section~\ref{sec:sr-pineau}). Deux
faits en expliquent une part : notre grille, qui trouve 35~\% en 2018, et les articles sur les modèles de
langue, qui décrivent rarement leurs références, si bien que tout le premier tirage tombe à 18~\% (8 sur
44).

\paragraph{F8, la mesure de performance}
La mesure est définie dans 78~\% des articles de 2026 contre 79~\% en 2018 (Y1). Quand des barres
d'erreur figurent, leur sens est défini dans 38 articles sur 44 en 2026 contre 25 sur 30 en 2018 (Y2).
La mesure est le plus souvent définie, les deux années, sans changement détectable. Y1 reste notre
critère le plus fragile : sur le premier tirage, deux lecteurs ne s'accordaient que sur 60~\% des
articles avant que ses règles soient précisées (section~\ref{sec:sr-accord}).

\paragraph{F9, le nombre d'essais}
Sur les 41 articles de RL expérimentaux du premier tirage (familles \file{llm} et \file{classic}), 20
ne disent pas combien d'entraînements indépendants soutiennent leurs résultats (X1). Quatre n'en ont
qu'un. La médiane des autres vaut cinq, la valeur la plus fréquente du tableau que Pineau montrait en
2018. Cette faille-là n'a pas été écoutée. Le chapitre~\ref{chap:06-experiences-resultats} mesure avec l'expérience E5
ce que cinq graines aléatoires permettent de trancher.

\paragraph{F10, les meilleurs essais}
La sélection des meilleurs essais se voit dans 14~\% des articles de 2026 contre 19~\% en 2018 (X3),
sans changement détectable. Dans le premier tirage, 2 articles sur 45 la laissent voir. Un article qui
garde ses meilleurs essais sans le dire ni le montrer échappe à notre lecture : ces parts sont des
minorants.

\paragraph{F12, l'incertitude et le test}
La faille se coupe en deux. Rapporter une incertitude est devenu courant : 88~\% des articles de 2026
contre 63~\% en 2018. Le type d'incertitude est nommé dans 41 des 44 articles de 2026 qui en rapportent
une, contre 25 sur 30 en 2018 (Y11). Nommer le type d'incertitude ne dit pas encore sur quoi elle est
calculée. Parmi les articles qui tracent une barre, la part de ceux qui la calculent sur des graines
aléatoires reste voisine d'un sur deux, 57~\% en 2018 et 48~\% en 2026 (X8, \cref{tab:x8}) : la leçon de la
faille F11 et de l'expérience B1 reste à entendre. Le test statistique n'a pas bougé de façon détectable : 12~\% des
articles de 2026 contre 6~\% en 2018 (Q8). Les outils apparus après 2018, comme l'IQM ou la bibliothèque
rliable, passent de 2~\% à 14~\% (Y12), sans passer la correction. Les articles tracent des barres, ils
testent rarement.

\paragraph{F13 et F14, les tâches et les implémentations}
La conclusion principale repose sur plusieurs tâches dans 41 articles sur 50 en 2026, contre 43 sur 48
en 2018 (Y3). Parmi les articles à plusieurs tâches, 38 sur 41 montrent au moins une tâche où leur
méthode ne gagne pas, contre 34 sur 43 en 2018 (Y4). L'origine du code des références est dite dans
41~\% des articles de 2026 contre 46~\% en 2018 (Y5). Aucun des trois critères ne change de façon
détectable. Les deux premières pratiques étaient déjà courantes en 2018, la troisième reste
minoritaire. Dans le premier tirage, 25 articles sur 45 seulement comptent plusieurs tâches : beaucoup
d'articles sur les modèles de langue s'entraînent sur un seul jeu de données.

\paragraph{F15, l'évaluation tenue à part}
À grille égale, le RL classique passe de 33~\% à 45~\% pour Q7 et de 23~\% à 34~\% pour le test de
généralisation, plus exigeant (Y8). Aucune des deux différences n'est détectable. La forte hausse apparente
contre les 10~\% de Pineau vient de la grille et de l'arrivée des articles sur les modèles de langue (sections~\ref{sec:sr-pineau} et~\ref{sec:sr-premier}).

\paragraph{F16, les bancs d'essai}
La version de chaque banc d'essai est donnée dans 26~\% des articles de 2026 contre 23~\% en 2018 (Y7),
sans changement détectable. Dans le premier tirage, 16 articles sur 27 de la famille \file{llm}
n'essaient qu'une seule famille de modèles de langue (X6).

\section{Ce qui résiste aux analyses de robustesse}\label{sec:sr-robustesse}

Le tableau~\ref{tab:robustesse} rassemble chaque analyse alternative et dit, pour chaque verdict, s'il la
traverse.

\gentable{robustesse}

\paragraph{La provenance des articles de 2026 (a1)}
Au sein de 2026, les 14 articles du premier tirage et les 36 de la file A3 ne diffèrent sur aucun des
quatre critères (Fisher, $p = 0{,}73$ pour Q5, 1 pour Y9 et X7, $0{,}17$ pour X2). Refaite avec chaque
provenance seule contre 2018, la comparaison garde son signe et un intervalle qui exclut zéro dans les
huit cas. Pour Q5 : $+47$ points $[+18\,;\,+65]$ avec le premier tirage, $+41$ $[+20\,;\,+57]$ avec la
file. Pour X2, l'intervalle de la file seule s'approche de zéro : $+21$ points $[+1\,;\,+38]$.

\paragraph{La composition par conférence (a3)}
Le test de Cochran-Mantel-Haenszel, stratifié par conférence, garde les quatre verdicts : $p = 5 \times
10^{-5}$ pour Q5, $0{,}001$ pour Y9, $0{,}005$ pour X7, $0{,}006$ pour X2. Il ne soutient pas Q6
($p = 0{,}06$) ni Y12 ($p = 0{,}15$).

\paragraph{La source du PDF (a2)}
Cette coupe oppose 45 articles de 2026 lus sur arXiv à 5 lus sur le site d'une conférence : elle ne peut
rien trancher et nous ne nous appuyons pas sur elle. La question qu'elle visait demande une mesure directe. Une
version arXiv est parfois plus longue ou plus récente que celle des actes, alors que les articles de 2018 ont
été lus dans les actes. Pour Q5, cette mesure a été faite.

\paragraph{Q5 relu sur la même version les deux années (écart 31)}
Le lien vers le code est le critère le plus exposé à l'asymétrie des versions : un auteur ajoute parfois
l'adresse de son dépôt dans une version arXiv postérieure à la conférence. L'écart 31, écrit avant tout
téléchargement, relit Q5 de trois façons qui traitent les deux années de la même manière :
\begin{itemize}
  \item \textbf{les deux années sur les actes} : les 45 articles de 2026 lus sur arXiv ont tous été retrouvés
  dans les actes ;
  \item \textbf{les deux années sur la dernière version arXiv} : 27 des 33 articles de 2018 qui valaient
  \file{no} y ont été retrouvés, les six autres gardent leur lecture ;
  \item \textbf{le code disponible à la conférence} : un \file{yes} dont tous les dépôts GitHub ont été créés
  après le dernier jour de la conférence devient \file{no}.
\end{itemize}
Un arbitre neuf a tranché les cas douteux à l'aveugle, années masquées et identifiants remplacés. Chaque lecture
reçoit aussi une borne où les cas incertains comptent \file{yes}. La borne de la dernière compte en plus
\file{no} les dépôts sans date. Le verdict tient dans les sept lectures, la lecture publiée, les trois relectures et leurs trois bornes
(\cref{tab:q5-versions,fig:q5-lectures}) : $+41$
points $[+21\,;\,+56]$ sur les actes, $+34$ $[+15\,;\,+51]$ sur la dernière version arXiv, $+37$
$[+17\,;\,+52]$ pour le code disponible à la conférence. La lecture la plus défavorable, la dernière version
arXiv avec les cas incertains comptés \file{yes}, donne encore $+30$ points $[+11\,;\,+47]$,
$p_{\text{Holm}} = 0{,}035$. L'asymétrie existe : en 2018, quatre articles n'ont un lien vers leur code que
dans leur dernière version arXiv (C19, C216, C22, C39). Elle touche quatre à six articles sur 48 et n'explique
qu'une petite part de la hausse. Un seul dépôt a été créé après sa conférence, celui de C219, un an après
NeurIPS 2018. Une réserve : l'arbitre compte C19 \file{yes} dans la lecture arXiv, alors que l'audit de l'écart 19 avait
visité sa page et jugé qu'elle ne menait qu'à une bibliothèque publiée après l'article. Sans C19, cette lecture
donne $+37$ points.

\gentable{q5-versions}

\fig{q5-lectures}{Q5, le code en lien, relu sur la même sorte de version les deux années : les actes, la
dernière version arXiv, le code disponible à la conférence, chacune avec sa borne. À gauche, les parts de 2018
et de 2026, avec les 55~\% de Pineau en pointillé. À droite, la différence et son intervalle de Newcombe à 95~\%.
La différence reste entre $+30$ et $+43$ points et passe la correction de Holm dans chaque lecture. Notre lecture de
2018 reste sous les 55~\% de Pineau dans toutes. Données : \file{study/q5\_versions.json},
\file{study/results\_final.json} (\cref{tab:q5-versions}).}

\paragraph{Un lecteur retiré à la fois (a4)}
Retirer tour à tour chacun des seize lecteurs principaux ne rend fragile aucun des quatre verdicts. La
borne basse la plus faible de l'intervalle de Newcombe vaut $+16$ points pour Q5 (sans H1 ou sans H3),
$+10$ pour Y9 (sans H2), $+7$ pour X7 (sans E, K2 ou K5) et $+4$ pour X2 (sans H4). Q6 devient fragile
sans le lecteur H2 ($+18$ points, $[-3\,;\,+36]$). Y12 est fragile dans 8 des 16 coupes.

\paragraph{Sans les treize articles lus en partie par un autre modèle}
Sans les treize articles finis par \file{claude-opus-5} (écart 29), l'échantillon de RL classique de
2026 passe à 48 articles et celui de 2018 reste à 50. Le code (Q5, $p_{\text{Holm}} = 0{,}0001$), la
sensibilité (Y9, $0{,}008$) et le calcul (X7, $0{,}04$) passent toujours la correction. L'incertitude
rapportée ne la passe plus : $+25$ points, $[+8\,;\,+41]$, $p_{\text{Holm}} = 0{,}07$. Q2 ne la passe
plus non plus ($p_{\text{Holm}} = 0{,}07$). Les tendances de la trajectoire ne changent pas
(section~\ref{sec:sr-trajectoire}).

\paragraph{La recherche contre les \enquote{non}}
Sur les sept critères vérifiés, les deux échantillons comptent 377 décisions finales \file{no}. Pour
94, les motifs fixés d'avance ne trouvent aucune occurrence dans le texte. Les 283 autres ont été
soumises à l'arbitre. L'arbitre en retourne 10 : 2 en 2018 (C19 pour Q5, C120 pour Y9) et 8 en 2026 (K058,
K063 et LR03 pour Q4, K002 pour Q7, ML06 pour Q8, K056, LR03 et LR09 pour X7). Il en juge 24 douteuses.
Refaite avec les décisions retournées, la comparaison principale garde tous ses verdicts
(\file{study/falsify\_no.json}). La différence du calcul chiffré grandit, à $+37$ points $[+17\,;\,+53]$.
Celle des hyperparamètres des références tombe à $-1$ point. Si l'on retourne aussi les 24 décisions
douteuses, l'incertitude rapportée et la sensibilité passent juste au-dessus du seuil
($p_{\text{Holm}} = 0{,}07$ et $0{,}054$). Cette seconde borne n'était pas annoncée dans l'écart 27.

\paragraph{Bilan de robustesse}
Le code et le calcul passent la correction dans tous nos contrôles. Le code la passe aussi quand les deux années
sont lues sur la même version. La sensibilité aux réglages la passe
dans tous sauf la borne des décisions douteuses, où elle s'arrête à $p_{\text{Holm}} = 0{,}054$.
L'incertitude rapportée échoue dans deux contrôles, sans les treize articles et avec la borne des
décisions douteuses, avec $p_{\text{Holm}} = 0{,}07$ chaque fois. Nous la gardons parmi les quatre
changements, avec cette réserve.

\paragraph{Le coût et le risque de chaque pratique}
Les trois évaluateurs du coût et du risque de chaque pratique (section~\ref{sec:sm-robustesse})
s'accordent entre eux : leurs notes se corrèlent de 0,78 à 0,91 pour le coût, de
0,90 à 0,93 pour le risque (Spearman, par paire). Leurs notes ne prédisent pas les différences observées. La
corrélation de Spearman vaut $+0{,}16$ pour le coût ($p = 0{,}52$ par permutation) et $-0{,}04$ pour le
risque ($p = 0{,}89$). Le balayage d'un hyperparamètre (Y9), noté coûteux (médiane 4), a progressé. Les
versions des bancs d'essai (Y7) et des dépendances (Y10), notées presque gratuites (médiane 1), n'ont pas
bougé. La lecture \enquote{écoutée là où cela coûtait peu} est retirée.

\section{Notre lecture de 2018 contre les chiffres de Pineau}\label{sec:sr-pineau}

\paragraph{La différence entre les deux lectures}
La comparaison secondaire de l'écart 12 oppose notre lecture de 2018 aux chiffres de Pineau, sur les
six critères Q3 à Q8 où elle en donne un comparable, avec une correction de Holm sur ces six critères
(\file{extension\_stats.json}, champ \file{ours2018\_vs\_pineau}).
\begin{itemize}
  \item \textbf{Q3, hyperparamètres de la méthode} : 55~\% (26 sur 47) chez nous, 90~\% chez Pineau.
  Différence de $-35$ points, $[-50\,;\,-17]$, $p_{\text{Holm}} = 0{,}001$.
  \item \textbf{Q7, évaluation tenue à part} : 33~\% (16 sur 48) chez nous, 10~\% chez Pineau. Différence de
  $+23$ points, $[+7\,;\,+39]$, $p_{\text{Holm}} = 0{,}03$.
  \item \textbf{Q4, hyperparamètres des références} : 35~\% contre 60~\%. Différence de $-25$ points,
  $[-42\,;\,-5]$, $p_{\text{Holm}} = 0{,}06$.
  \item \textbf{Q5, code en lien} : 31~\% contre 55~\%. Différence de $-25$ points, $[-42\,;\,-5]$,
  $p_{\text{Holm}} = 0{,}06$.
  \item \textbf{Q6 et Q8} : 26~\% contre 20~\% et 6~\% contre 5~\%, sans différence détectable.
\end{itemize}

\gentable{sondage-pineau-contexte}

Les 12 premiers articles de la calibration donnaient déjà le même signe : Q3 à 6 sur 10, Q7 à 4 sur
11 (\file{results\_final.json}). Notre grille est plus sévère que celle de Pineau sur ce que l'article
donne, plus large sur ce qui compte comme tenu à part.

\paragraph{Pourquoi notre lecture diffère}
Le tirage de Pineau et sa façon de trancher les cas douteux restent inconnus : la diapositive annonce
50 articles de RL de 2018 publiés à NeurIPS, ICML et ICLR, sans dire comment ils ont été choisis ni
comment chaque critère a été décidé. Nous ne pouvons donc pas attribuer la différence à une cause unique.
Trois mécanismes le rendent plausible, sans qu'aucun soit démontré.
\begin{itemize}
  \item \textbf{Des règles écrites plus strictes sur ce qui est donné.} Nos règles de Q3 et Q4 exigent
  les hyperparamètres de toutes les méthodes principales. Des \enquote{valeurs par défaut} sans liste
  ne comptent que si la version de la bibliothèque est donnée. Pour Q5, \enquote{sera publié} et une page
  de projet qui ne mène pas au code valent \file{no}. Une lecture qui accepte un tableau partiel ou une
  promesse de publication donne des parts plus hautes.
  \item \textbf{Une règle plus large sur ce qui est tenu à part.} Notre Q7 compte des graines
  d'environnement non vues à l'entraînement. Pineau pensait sans doute à une séparation plus nette,
  puisque sa feuille de route de NeurIPS demande de séparer les tâches et de passer à des simulateurs
  plus riches.
  \item \textbf{Un autre tirage.} Notre échantillon de 2018 est tiré au hasard parmi les 131 articles
  dont le titre correspond à notre expression. Celui de Pineau peut contenir des articles que notre
  filtre ne voit pas, ou en exclure.
\end{itemize}
Un quatrième mécanisme a été mesuré pour Q5 : la version lue. Relue sur la dernière version arXiv, notre
lecture de 2018 passe de 31~\% à 40~\% (19 sur 48), 44~\% en comptant les cas incertains (21 sur 48,
\cref{tab:q5-versions}). La version explique une part de la différence avec les 55~\% de Pineau, pas toute la différence.
Après correction, les différences de Q4 et de Q5 restent compatibles avec le hasard ($p_{\text{Holm}} =
0{,}06$). Ceux de Q3 et de Q7 ne le sont pas.

\paragraph{La règle de verdict qui en découle}
Comparer nos articles de 2026 aux chiffres de Pineau mêlerait le changement des pratiques et celui de
la grille. L'écart 20 en tire la règle de verdict (section~\ref{sec:sm-verdict}). Les chiffres de Pineau restent
montrés comme contexte, à la section~\ref{sec:sr-premier}.

\paragraph{Ce qu'aurait donné la règle d'origine}
L'écart 30 (c), écrit avant le calcul du dépôt mais après une esquisse de la relecture technique
(section~\ref{sec:ecarts-audit}), applique aux mêmes données la règle d'origine de l'écart 7, que l'écart
20 a remplacée : chaque critère de Pineau était jugé contre ses chiffres, puis contre notre lecture de
2018 (section~\ref{sec:sm-verdict}, \cref{tab:regle-origine}). Dans le RL classique de 2026, aucun critère n'aurait été \enquote{changé et
confirmé}. Q4, Q6 et Q7 auraient \enquote{changé} contre Pineau, sans confirmation contre notre 2018. Q5, notre
verdict le plus fort, aurait été \enquote{ambigu} : il passe la correction contre notre lecture de 2018, pas contre
les chiffres de Pineau ($+18$ points, $p_{\text{Holm}} = 0{,}19$). Dans le premier tirage, grands modèles compris,
seul Q7 aurait été \enquote{changé et confirmé}, porté par les articles sur les modèles de langue
(section~\ref{sec:sr-premier}). La règle d'origine ne dit rien de X2, X7 et Y9 : Pineau ne donne pas de chiffre
pour eux. Ce tableau ne renverse pas la décision de l'écart 20, qui repose sur la différence de calibration mesurée
pour Q3, 55~\% chez nous contre 90~\% chez Pineau. Il rend cette décision vérifiable.

\gentable{regle-origine}

\section{Le premier tirage et les modèles de langue}\label{sec:sr-premier}

\paragraph{Le premier tirage contre Pineau}
La figure~\ref{fig:remesure-dumbbell} compare les 45 articles expérimentaux du premier tirage, grands
modèles de langue compris, aux chiffres de Pineau. La figure~\ref{fig:ecarts-pineau} dessine les différences avec
leurs intervalles de Newcombe, la figure~\ref{fig:remesure-ic} les parts de 2026 avec leurs intervalles de
Wilson. Trois intervalles s'écartent nettement de zéro. Les hyperparamètres de la méthode baissent, de
90~\% à 67~\% ($-23$ points, $[-39\,;\,-7]$, $p_{\text{Holm}} = 0{,}02$). Ceux des références baissent,
de 60~\% à 18~\% ($-42$ points, $[-57\,;\,-22]$, $p_{\text{Holm}} = 0{,}0002$). L'évaluation tenue à
part monte, de 10~\% à 70~\% ($+60$ points, $[+42\,;\,+73]$). Le code passe de 55~\% à 67~\% sans différence
détectable. Le test statistique passe de 5~\% à 4~\%. Deux causes s'y cumulent, le changement des
pratiques et la différence des grilles. Ces différences servent de contexte, jamais de verdict.

\fig{remesure-dumbbell}{La grille de Pineau, appliquée en 2018 et en 2026, à titre de contexte. Chaque
ligne est un critère de lecture. Le point de 2018 donne la part des 50 articles lus par Pineau qui
satisfont le critère. Le point de 2026 donne la part dans les 45 articles expérimentaux de notre
premier tirage d'ICLR 2026, d'ICML 2026 et de NeurIPS 2025, grands modèles de langue compris. Les deux
points ne viennent pas de la même grille : notre lecture de 2018 trouve 55~\% pour les hyperparamètres
de la méthode, là où Pineau trouvait 90~\%. Le test statistique, absent du dessin, passe de 5~\% à
4~\%. Données : \file{study/audit\_stats.json}.}

\fig{ecarts-pineau}{Différence entre la part des 45 articles expérimentaux de notre premier tirage de 2026 et
celle des 50 articles de 2018 lus par Pineau, critère par critère, en points de pourcentage. Le point
est la différence observée, la barre son intervalle de Newcombe à 95~\%, qui tient compte de l'incertitude des
deux échantillons. Au-dessus de zéro, la pratique est plus fréquente en 2026. La différence mêle le
changement des pratiques et la différence entre notre grille et celle de Pineau : il sert de contexte,
les verdicts viennent de la comparaison à grille égale. Données : \file{study/audit\_stats.json}.}

\fig{remesure-ic}{Part des 45 articles expérimentaux du premier tirage de 2026 qui satisfont chaque
critère, avec son intervalle de Wilson à 95~\%. Le tirage mêle RL classique et articles sur les grands
modèles de langue. Avec un échantillon de cette taille, une différence de moins de quinze points avec 2018 se
confond avec le hasard du tirage. Données : \file{study/results\_final.json}.}

\paragraph{Les articles sur les modèles de langue}
Le tableau~\ref{tab:sondage-lrl} oppose, dans le premier tirage, les 27 articles de la famille \file{llm} aux
14 articles de RL classique, que plusieurs critères séparent nettement. L'évaluation tenue à part est presque
universelle dans la famille \file{llm} (26 sur 27) et minoritaire dans le RL classique du premier tirage
(5 sur 13) : les recueils de problèmes de mathématiques séparent par construction les problèmes
d'entraînement et de test. L'incertitude est rare dans la famille \file{llm} (4 sur 27) et universelle
dans le RL classique du premier tirage (14 sur 14). Les hyperparamètres des références sont tous donnés
dans 5 articles \file{llm} sur 27. Seize articles \file{llm} sur 27 n'essaient qu'une famille de modèles
de base. La hausse apparente de l'évaluation tenue à part contre Pineau vient presque entière de ces
articles. La baisse apparente de l'incertitude dans le premier tirage aussi.

\gentable{sondage-lrl}

\paragraph{Un changement de population}
La famille \file{llm} n'existait pas en 2018. Une différence entre le premier tirage de 2026 et 2018 qui
passe par ces articles mesure d'abord l'arrivée d'une population nouvelle, avec ses propres habitudes.
Il ne dit rien du changement des pratiques au sein d'une population restée comparable. Les verdicts de
ce rapport valent donc pour le RL classique seul.

\section{La trajectoire de 2018 à 2026}\label{sec:sr-trajectoire}

\paragraph{Les parts année par année}
Le tableau~\ref{tab:trajectoire} et la figure~\ref{fig:trajectoire} donnent, pour chaque pratique, la part des
articles de RL classique qui la suivent en 2018 (48 articles), en 2019, 2020, 2022 et 2024 (20 articles
par année, ICML et NeurIPS) et en 2026 (50 articles). Le code en lien passe de 31~\% en 2018 à 60~\%,
50~\%, 70~\% et 75~\% de 2019 à 2024, puis 74~\% en 2026. La sensibilité à un hyperparamètre vaut 44~\%,
puis 35~\%, 45~\%, 60~\%, 85~\% et 80~\%. Le calcul chiffré vaut 27~\%, puis 30~\%, 40~\%, 40~\%, 55~\%
et 58~\%. Une incertitude rapportée vaut 63~\%, puis 80~\%, 85~\%, 95~\%, 100~\% et 88~\%. Le test
statistique vaut 6~\%, puis 5~\%, 10~\%, 10~\%, 25~\% et 12~\%. Avec 20 articles, une part a une
incertitude d'environ vingt points de part et d'autre.

\gentable{trajectoire}

\fig{trajectoire}{Part des articles de RL classique qui suivent chaque pratique, année par année : 48
articles de 2018 (ceux qui comportent des expériences), 20 par année pour 2019, 2020, 2022 et 2024, tirés d'ICML et de
NeurIPS, 50 de 2026. La ligne verticale en pointillé marque la liste de contrôle de NeurIPS (2019). Le
code en lien, la sensibilité à un hyperparamètre, le calcul chiffré et l'incertitude rapportée montent
régulièrement. Le test statistique, en gris, ne montre aucune tendance. Avec 20 articles, une part a
une incertitude d'environ vingt points de part et d'autre. Données : \file{study/trajectory.json}.}

\paragraph{Les tendances}
Le test de Cochran-Armitage, avec les années pour scores, donne une tendance nette pour les quatre
pratiques après correction de Holm sur les cinq critères testés : $z = 4{,}29$ pour le code
($p_{\text{Holm}} = 7 \times 10^{-5}$), $z = 4{,}78$ pour la sensibilité ($9 \times 10^{-6}$),
$z = 3{,}44$ pour le calcul ($0{,}002$), $z = 3{,}39$ pour l'incertitude ($0{,}002$). Le test
statistique ne montre aucune tendance ($z = 1{,}53$, $p = 0{,}13$).

\paragraph{Aucun saut juste après 2019}
De 2018 aux années 2019 et 2020 réunies, aucune différence ne passe la correction. Le code monte de $+24$ points, $[+3\,;\,+42]$,
avec $p = 0{,}03$ avant correction et $p_{\text{Holm}} = 0{,}16$ après. L'incertitude monte de $+20$
points, $[+1\,;\,+37]$, $p_{\text{Holm}} = 0{,}23$. La sensibilité ne monte pas ($-4$ points). Des
années 2019 et 2020 à 2022 et 2024, seule la sensibilité monte nettement, de 40~\% à 73~\%, avec un intervalle de
$+11$ à $+50$ points et $p_{\text{Holm}} = 0{,}03$. Ces séries décrivent une hausse lente, commencée avant
la liste de contrôle ou avec elle. Elles ne disent pas ce qui l'a causée. \textcite{coakley2026shift}
voient de même la part du code publié monter avant les listes de contrôle, sur cinq congrès
d'intelligence artificielle.

\paragraph{Sans les treize articles}
Sans les articles finis par \file{claude-opus-5}, les quatre tendances passent toujours la correction
($p_{\text{Holm}} < 0{,}004$) et le test statistique n'en montre toujours aucune. La montée tardive de
la sensibilité garde $p_{\text{Holm}} = 0{,}02$ (\file{trajectory\_sans\_opus5.json}).

\paragraph{Le pic du test statistique en 2024}
Le test statistique atteint 25~\% des articles en 2024 (5 sur 20), puis retombe à 12~\% en 2026 (6 sur
50). Avec 20 articles, ce pic reste dans la marge d'incertitude. Nous ne le lisons pas comme un fait.

\section{L'accord entre lecteurs}\label{sec:sr-accord}

Le tableau~\ref{tab:accord} et la figure~\ref{fig:accord-lecteurs} donnent l'accord des doubles lectures à
l'aveugle, critère par critère, sur les lectures brutes avant harmonisation.

\gentable{accord}

\fig{accord-lecteurs}{Accord entre les doubles lectures à l'aveugle, critère par critère, sur les
lectures brutes, avant harmonisation : part de décisions identiques et kappa de Cohen quand les deux
lecteurs varient. Quatre paires de lectures de dix articles chacune : G (grille v1) et G2 (grille v2)
sur le premier tirage de 2026, R2 sur les articles ajoutés de 2018 et de 2026, R3 sur la trajectoire de
2019 à 2024. Le seuil de 70~\% d'accord, sous lequel un critère est dit fragile, est marqué. Données :
\file{study/agree\_v2.json}, \file{study/audit\_stats.json}.}

\paragraph{Trente articles relus, avec leurs intervalles}
Les quatre paires de doubles lectures réunissent 30 articles distincts. G et G2 relisent les mêmes dix articles
de 2026, l'un avec la grille v1, l'autre avec la grille v2. R2 relit dix articles ajoutés, cinq de 2018 et cinq
de 2026, R3 dix articles de la trajectoire. L'écart 30 (a), écrit avant ce calcul, réunit ces lectures et donne
à chaque part d'accord son intervalle de Wilson, à chaque kappa un intervalle bootstrap percentile sur les
articles (\cref{tab:accord-intervalles}). Sur les quatre critères des verdicts, les 120 décisions doublées sont
identiques : 30 sur 30 pour chaque critère, avec un intervalle de Wilson de 89~\% à 100~\%. Trois critères
gardent un kappa sous 0,70 : la mesure définie (Y1), l'origine du code des références (Y5) et les exclusions
dites (Y6). Une paire seule renseigne mal : sur dix articles, 9 accords ont pour intervalle [60\,;\,98]~\%. Le
kappa de 0,80 de G sur Q3 a pour intervalle [0,31\,;\,1]. L'accord sur des articles de 2018 repose sur les cinq
articles de R2 : 5 accords sur 5 ont pour intervalle de Wilson [57\,;\,100]~\%.

\gentable{accord-intervalles}

\paragraph{La grille v1 sur le premier tirage}
Le second lecteur G et les lecteurs principaux s'accordent sur 127 décisions sur 130 (98~\%), soit 13
critères sur 10 articles. Les trois désaccords touchent Q3, Q7 et X4, avec des kappas de 0,80, 0,78 et
0,74. Le kappa vaut 1 pour les autres critères où les deux lecteurs varient.

\paragraph{La grille v2 sur le premier tirage}
Sur les mêmes dix articles, G2 et les lecteurs principaux s'accordent complètement sur 9 critères sur 12.
L'accord tombe à 90~\% sur les exclusions (Y6, kappa 0,80), à 80~\% sur l'implémentation des
références (Y5, kappa 0,52) et à 60~\% sur la mesure de performance (Y1, kappa 0,31). Y1 est le seul
critère fragile de toutes les doubles lectures. Ses règles ont été précisées ensuite par les écarts 17
et 19.

\paragraph{Les articles ajoutés}
R2 relit dix articles tirés parmi les ajouts de 2018 et de 2026 avec les deux grilles. L'accord atteint
au moins 80~\% sur chacun des 23 critères et il est complet sur 17 d'entre eux. Les accords les plus bas
sont ceux de Y4 (80~\%, kappa 0,52) et de Y6 (80~\%, kappa 0,41).

\paragraph{La trajectoire}
R3 relit dix articles de A4, avec le même bilan : aucun des 23 critères ne descend sous 80~\% d'accord
et 17 restent complets, dont les cinq critères testés de la trajectoire (Q5, Y9, X7, X2, Q8). Les kappas les
plus bas sont ceux de Q7 et de Y8 (80~\%, kappa 0,41).

\paragraph{La stabilité d'un seul modèle}
Tous les lecteurs sont des copies du même modèle, hormis les deux qui ont achevé treize articles sur
\file{claude-opus-5} (section~\ref{sec:sm-lecteurs}). L'accord des doubles lectures mesure donc la
stabilité de ce modèle (section~\ref{sec:sm-doubles}). Un humain aurait pu
lire autrement les mêmes passages. Les citations et les pages de \file{study/codes/} et
\file{study/codes\_v2/} permettent à chacun de vérifier chaque décision. Un lecteur d'une autre famille de
modèles, \file{gemma3:12b}, relit les 98 articles des deux échantillons sur les critères des verdicts (écart 33,
écrit avant cette lecture). Sa lecture n'était pas finie quand ce rapport a été compilé et n'y figure pas.

\paragraph{L'année et les lecteurs démêlés}\label{sec:sr-annees}
Dans la comparaison principale, les lecteurs de 2018 ne sont pas ceux de 2026
(section~\ref{sec:sm-annees}). L'écart 36 fait relire les 98 articles expérimentaux de cette
comparaison par huit agents neufs, T1 à T8, qui lisent chacun six articles de 2018 et six ou sept de
2026 (section~\ref{sec:sm-doubles}). Avec ces seules lectures, la part de \enquote{oui} passe de 31~\%
à 72~\% pour le code en lien (Q5), de 42~\% à 78~\% pour la sensibilité à un hyperparamètre (Y9), de
31~\% à 66~\% pour le calcul chiffré (X7) et de 69~\% à 86~\% pour l'incertitude rapportée (X2)
(\cref{tab:r36}).

\gentable{r36}

Le modèle logistique mixte donne à chaque lecteur son propre penchant pour le \enquote{oui}, puis
mesure l'effet de l'année à lecteur égal. L'écart type $s$ de l'effet lecteur vaut 0 pour Q5 et Y9, 0,24 pour
X7 et 0,44 pour X2 : aucun ne se distingue de zéro ($p \geq 0{,}25$). Le penchant pour le
\enquote{oui} varie donc peu d'un lecteur à l'autre. Après la correction de Holm sur les quatre
critères, Q5, Y9 et X7 passent les deux tests fixés d'avance, celui de Fisher sur la différence et
celui de l'effet de l'année dans le modèle. Leurs verdicts tiennent. Pour X2, la différence vaut
$+17$ points, $[+0{,}6\,;\,+33{,}0]$. L'effet de l'année passe la correction
($p_{\text{Holm}} = 0{,}037$), mais le test de Fisher la manque de peu ($p_{\text{Holm}} = 0{,}053$) :
la règle dit ce verdict \enquote{même sens, non établi}. X2 était déjà le verdict le plus fragile
(section~\ref{sec:sr-robustesse}). Chaque intervalle de Newcombe chevauche celui de la comparaison
principale. Dans la lecture la moins favorable à Q5, où toute page de projet compte \enquote{oui},
sa hausse passe encore la correction de Holm, de 42~\% à 74~\% ($p_{\text{Holm}} = 0{,}004$).

\paragraph{Les lectures neuves rejoignent les codes d'origine}
Chaque article a ainsi deux lectures complètes : celle d'un lecteur neuf et le code final d'origine,
harmonisé. Leur accord va de 94~\% à 99~\% selon le critère, de 90~\% à 100~\% selon le critère et
l'année (\cref{tab:r36-accord}). Aucun critère ne montre d'accord nettement différent entre 2018 et
2026 : le plus petit $p$ de Fisher vaut 0,11. Pour X7, les six désaccords vont tous vers le
\enquote{oui}, deux en 2018 et quatre en 2026. Huit \enquote{oui} de X7 reposent sur des durées
mesurées (temps d'exécution, d'inférence ou d'entraînement) plutôt que sur un calcul total déclaré.
L'écart 19 laisse ce cas limite au jugement du lecteur. Trois de ces articles sont de 2018 et cinq de
2026. Six des huit valent déjà \enquote{oui} dans le code d'origine. Le
chapitre~\ref{chap:10-limites} dit ce que cette relecture ne sépare pas.

\gentable{r36-accord}

\section{Ce que les articles déclarent}\label{sec:sr-declarations}

\paragraph{La signification statistique}
La question de la liste de contrôle de NeurIPS sur la signification statistique se contente de barres
d'erreur \enquote{ou d'une autre information appropriée}. Sur les 30 articles de NeurIPS 2025 lus, 18
y répondent \enquote{Yes}. Aucun des 18 n'applique de test : la liste ne demande pas le test et les
articles ne le font pas. Un seul des 18 ne rapporte pas non plus d'incertitude (\file{study/checklist.json}).

\paragraph{Le code, les réglages et le calcul}
Sur le code, 22 articles répondent \enquote{Yes}, dont 5 sans lien vers leur code dans la version que
nous avons lue, soit 23~\% $[10\,;\,43]$. Quatre n'ont pas de lien même dans la version des actes. Trois
écrivent dans leur propre justification que le code sera publié plus tard. Sur les réglages, 22 articles
répondent \enquote{Yes}, dont 6 sans donner tous les hyperparamètres de leur méthode selon notre règle.
Sur le calcul, 21 articles répondent \enquote{Yes}, dont 6 sans chiffrer de temps de calcul, que la
question demande pourtant. Un article n'a pas de liste lisible dans les actes. Une case cochée ne
garantit pas la pratique qu'elle nomme.

\section{La sensibilité au filtre de titre}\label{sec:sr-filtre}

\paragraph{Combien d'articles de RL le filtre manque}
Sur les 400 articles tirés hors du filtre de titre, 25 ont un résumé qui correspond à l'expression
régulière : 6,2~\%, avec un intervalle de Wilson de 4,3~\% à 9,1~\% (\file{study/title\_filter.json}).
Rapportée à la population hors filtre, cette part donne environ \num{1037} articles de RL manqués, de
708 à \num{1503}, soit 45~\% de l'ensemble des articles de RL selon ce critère (de 36~\% à 55~\%). C'est
un majorant : la règle, mécanique, compte aussi des faux positifs, comme un article où \enquote{RL}
désigne l'apprentissage de représentations. La plupart des articles trouvés traitent des grands
modèles de langue.

\paragraph{Les vingt articles de RL caché}
Lus avec les deux grilles, les 20 premiers articles de RL caché ressemblent au premier tirage de 2026 :
le code en lien dans 60~\% d'entre eux (67~\% dans le tirage), une incertitude dans 50~\% (47~\%).
Treize traitent des grands modèles de langue. Les 7 articles de RL classique sont trop peu nombreux pour
éprouver nos verdicts. Le filtre rate beaucoup d'articles, surtout sur les modèles de langue, sans signe
que ces articles suivent d'autres pratiques.

\section{Synthèse}\label{sec:sr-synthese}

À grille égale, dans le RL classique, quatre pratiques ont changé entre 2018 et 2026 : le lien vers le
code, la sensibilité à un hyperparamètre montrée, le calcul chiffré, une incertitude rapportée. Le code
et le calcul résistent à toutes les coupes, la sensibilité à presque toutes. L'incertitude rapportée
échoue dans deux contrôles, avec $p_{\text{Holm}} = 0{,}07$. Huit lecteurs neufs, qui lisent chacun les
deux années, retrouvent les hausses du code, de la sensibilité et du calcul. Celle de l'incertitude
garde son sens sans être établie. Le lien vers le code garde aussi son verdict quand les deux années
sont lues sur la même version. Les quatre pratiques montent lentement de 2018 à 2024, sans saut juste
après la liste de contrôle de 2019.

Le test statistique, les références décrites et réglées avec soin, l'évaluation tenue à part, les
versions des bancs d'essai et des bibliothèques n'ont pas bougé de façon détectable. Le réglage des
références ne montre aucun signe de progrès. Pour le reste, le sondage ne peut pas trancher, puisqu'il
ne voyait que des hausses d'environ trente points. La barre d'erreur s'est répandue sans que la part
calculée sur des graines aléatoires bouge de façon détectable (X8). Les changements nets concernent ce
qui accompagne le résultat : le lien vers le code, le calcul chiffré, l'effet d'un réglage, la barre
d'erreur.

Notre lecture de 2018 diffère de celle de Pineau sur les hyperparamètres de la méthode et sur
l'évaluation tenue à part, d'où la règle de verdict à grille égale. La règle de verdict d'origine,
appliquée aux mêmes données, n'aurait confirmé aucun changement dans le RL classique. Le
chapitre~\ref{chap:10-limites} rassemble les limites de ces résultats. Le
chapitre~\ref{chap:11-conclusion} les met en regard des expériences.
